% =============================================================================
% ForensiQ — ADR-0001 — Base de Dados e Alojamento: PostgreSQL + Neon
% UC 21184 · Projeto de Engenharia Informática · UAb
% =============================================================================
\documentclass[a4paper,11pt]{article}

\usepackage[utf8]{inputenc}
\usepackage[T1]{fontenc}
\usepackage[portuguese]{babel}
\usepackage[top=2.5cm, bottom=2.5cm, left=3cm, right=2.5cm]{geometry}
\usepackage{booktabs}
\usepackage{longtable}
\usepackage{array}
\usepackage{parskip}
\usepackage{hyperref}
\usepackage{enumitem}

\hypersetup{
    colorlinks=true,
    linkcolor=black,
    urlcolor=blue,
    pdfauthor={João Rodrigues},
    pdftitle={ForensiQ — ADR-0001 Base de Dados e Alojamento}
}

\newcolumntype{L}[1]{>{\raggedright\arraybackslash}p{#1}}

% =============================================================================
\begin{document}
% =============================================================================

\begin{center}
    {\LARGE\bfseries ADR-0001 --- Base de Dados e Alojamento:\\[0.3em]
    PostgreSQL + Neon}\\[1.2em]
    \begin{tabular}{ll}
        \textbf{Data:}       & 2026-03-23 \\
        \textbf{Estado:}     & Aceite \\
        \textbf{Decisores:}  & João Rodrigues \\
    \end{tabular}
\end{center}

\vspace{0.5em}
\hrule
\vspace{1.5em}

% =============================================================================
\section{Contexto}
% =============================================================================

O ForensiQ é uma plataforma web de gestão de prova digital para \emph{first
responders} da PSP, construída com Django + Django REST Framework. A aplicação
requer:

\begin{itemize}[nosep, leftmargin=*]
    \item \textbf{Integridade da cadeia de custódia} --- toda a alteração ou
    eliminação de registos de prova deve ser auditável e imutável, com
    conformidade à norma ISO/IEC 27037.

    \item \textbf{Hashing criptográfico} --- os ficheiros de prova são
    identificados por \emph{digest} SHA-256, calculado e verificado ao nível
    da base de dados.

    \item \textbf{Controlo de acesso granular} --- perfis distintos
    (administrador, investigador, \emph{first responder}, perito) com
    permissões diferenciadas a nível de linha.

    \item \textbf{Credibilidade forense e legal} --- a tecnologia escolhida
    deve ser reconhecida em contexto judicial e académico.

    \item \textbf{Custo controlado} --- o projecto é académico, com orçamento
    próximo de zero durante o desenvolvimento, mas com necessidade de ambiente
    funcional e público para a defesa pública (julho de 2026).
\end{itemize}

A escolha da base de dados e do fornecedor de alojamento são decisões
interligadas: o motor de base de dados determina as funcionalidades
disponíveis, e o fornecedor determina a sustentabilidade operacional durante
e após o período de desenvolvimento.

% =============================================================================
\section{Decisão}
% =============================================================================

\textbf{Decidimos utilizar PostgreSQL como motor de base de dados, alojado na
plataforma Neon (\href{https://neon.tech}{neon.tech}).}

Durante o desenvolvimento, será utilizado o plano gratuito do Neon
(\emph{Free Tier}). Antes da defesa pública (antecipadamente ao período
6--10 de julho de 2026), o projecto será migrado para o plano \emph{Neon
Launch} (5~USD/mês), de forma a garantir disponibilidade e desempenho
adequados para a demonstração.

% =============================================================================
\section{Alternativas consideradas}
% =============================================================================

\begin{longtable}{L{3.5cm} L{10.5cm}}
    \toprule
    \textbf{Alternativa} & \textbf{Razão de rejeição} \\
    \midrule
    \endhead
    \textbf{Supabase}
        & O plano gratuito pausa a base de dados ao fim de 7 dias de
          inactividade. O plano pago começa nos 25~USD/mês ---
          desproporcionado para um projecto académico. \\[0.5em]
    \textbf{Render (PostgreSQL managed)}
        & O plano gratuito expira ao fim de 30 dias, obrigando a recriar a
          instância repetidamente durante o desenvolvimento. \\[0.5em]
    \textbf{Aiven}
        & Opção tecnicamente válida, com suporte completo a PostgreSQL.
          Rejeitada por apresentar menor documentação de integração com
          Django e plano gratuito mais restritivo. Pode ser reavaliada em
          contexto de produção futura. \\[0.5em]
    \textbf{ElephantSQL}
        & Encerrado em janeiro de 2025. Não disponível. \\[0.5em]
    \textbf{Railway}
        & Não oferece \emph{tier} permanente sem custos; requer cartão de
          crédito e cobra após período de créditos iniciais. \\[0.5em]
    \textbf{SQLite}
        & Adequado apenas para desenvolvimento local. Não suporta
          \texttt{BEFORE UPDATE/DELETE} \emph{triggers}, \texttt{pgcrypto},
          nem acesso concorrente multi-utilizador. Incompatível com os
          requisitos forenses do projecto. \\
    \bottomrule
\end{longtable}

% =============================================================================
\section{Consequências}
% =============================================================================

\textbf{Positivas:}

\begin{itemize}[nosep, leftmargin=*]
    \item Os \emph{triggers} \texttt{BEFORE UPDATE} e \texttt{BEFORE DELETE}
    do PostgreSQL permitem implementar imutabilidade da cadeia de custódia
    directamente na base de dados, independentemente da camada aplicacional
    --- conformidade com ISO/IEC 27037.

    \item A extensão \texttt{pgcrypto} suporta o cálculo e verificação de
    \emph{digests} SHA-256 ao nível da base de dados, reforçando a
    integridade da prova mesmo em caso de falha da aplicação.

    \item O RBAC nativo do PostgreSQL complementa o sistema de permissões
    do Django, permitindo defesa em profundidade.

    \item O modelo ACID completo garante consistência transaccional em
    operações críticas (registo de prova, transferência de custódia).

    \item O Neon oferece dados que nunca expiram no plano gratuito; o
    \emph{auto-suspend} é transparente para o utilizador (retoma em
    milissegundos).

    \item O plano \emph{Neon Launch} (5~USD/mês) é suficiente para a demo
    pública e está dentro do orçamento do projecto.

    \item PostgreSQL é amplamente reconhecido em contextos forenses, judiciais
    e académicos, reforçando a credibilidade do projecto na defesa pública.
\end{itemize}

\textbf{Negativas / trade-offs:}

\begin{itemize}[nosep, leftmargin=*]
    \item O \emph{auto-suspend} do Neon no plano gratuito pode introduzir uma
    latência de arranque (\emph{cold start}) na primeira ligação após período
    de inactividade. Este comportamento é aceitável em contexto de
    desenvolvimento; será eliminado com a migração para o plano \emph{Launch}
    antes da defesa.

    \item A dependência de um fornecedor externo (Neon) introduz risco de
    disponibilidade. Mitigação: os \emph{backups} locais e o facto de o
    projecto ser académico tornam este risco aceitável.

    \item A configuração de \texttt{pgcrypto} e dos \emph{triggers} de
    custódia requer atenção na gestão de migrações Django
    (\texttt{RunSQL} em vez de migrações automáticas para esses elementos).
\end{itemize}

\vspace{2em}
\hrule
\vspace{0.5em}
{\small Para criar um novo ADR: copiar \texttt{ADR-000-template.tex},
incrementar o número, preencher e actualizar o estado.}

\end{document}
